\documentclass[aspectratio=169]{beamer}

% Paquetes esenciales
\usepackage[utf8]{inputenc}
\usepackage[spanish]{babel}
\usepackage{amsmath, amssymb}
\usepackage{graphicx}
\usepackage{booktabs}
\usepackage{tikz}
\usepackage{circuitikz} % Para el ejercicio diagnóstico
\usepackage{hyperref}

% Configuración del tema (Estilo profesional y limpio)
\usetheme{Boadilla}
\usecolortheme{seahorse}
\setbeamertemplate{navigation symbols}{}
\setbeamertemplate{footline}[frame number]

% Metadatos de la presentación
\title[Circuitos Electrónicos I]{Fundamentos Matriciales e Ingeniería Basada en Evidencia}
\subtitle{Sesión 01: Introducción y Puente Matricial}
\author[M.Sc. Ing. Bernardo Quiroga T.]{M.Sc. Ing. Bernardo Quiroga Turdera}
\institute[UCB Tarija]{
    Carrera de Ingeniería Mecatrónica\\
    Universidad Católica Boliviana ``San Pablo''\\
    Sede Tarija
}
\date{3 de Agosto de 2026}

\begin{document}

% ---------------------------------------------------
% Diapositiva de Título
% ---------------------------------------------------
\begin{frame}
    \titlepage
\end{frame}

% ---------------------------------------------------
% Filosofía del Curso
% ---------------------------------------------------
\begin{frame}{Filosofía del Curso: Estándar CDIO}
    \textbf{¿Qué diferencia a un técnico de un ingeniero mecatrónico o biomédico?}
    \vspace{0.5cm}
    
    En esta asignatura operaremos bajo el estándar \textbf{CDIO} (Concebir, Diseñar, Implementar, Operar) y los principios de la \textbf{Ingeniería Basada en Evidencia}.
    
    \vspace{0.5cm}
    \begin{itemize}
        \item \textbf{Cero pasividad:} No memorizaremos ecuaciones abstractas sin contexto.
        \item \textbf{Modelado Físico:} Traduciremos topologías eléctricas a modelos matemáticos.
        \item \textbf{Validación computacional:} Usaremos simuladores (gemelos digitales) antes de encender una fuente real.
        \item \textbf{Control de Calidad:} Todo hardware construido debe coincidir con la predicción matemática con un margen de error $\epsilon_r \le 5\%$.
    \end{itemize}
\end{frame}

% ---------------------------------------------------
% Proyecto Final Integrador (PFI)
% ---------------------------------------------------
\begin{frame}{Proyecto Final Integrador (PFI)}
    \begin{block}{Sistema Inteligente de Acondicionamiento de Señal}
        Diseño, simulación y validación de hardware para adquisición de bioseñales y sensores industriales.
    \end{block}
    
    \vspace{0.3cm}
    \textbf{Aplicaciones por Perfil Profesional:}
    \begin{itemize}
        \item \textbf{Ingeniería Biomédica:} Captura y filtrado analógico de señales bioeléctricas (EMG/ECG) con supresión de ruido de red (50 Hz).
        \item \textbf{Ingeniería Mecatrónica:} Acondicionamiento de celdas de carga y galgas extensiométricas para monitoreo de torque en lazo cerrado.
    \end{itemize}
    
    \vspace{0.3cm}
    \emph{Nota: Cada Elemento de Competencia del semestre aporta un bloque funcional a este sistema integrador.}
\end{frame}

% ---------------------------------------------------
% El Puente Matricial
% ---------------------------------------------------
\begin{frame}{El Pacto Pedagógico: El Puente Matricial}
    Sabemos que la formación matemática previa en álgebra tiene matices distintos en el aula. \textbf{Nadie está en desventaja.}
    
    \vspace{0.5cm}
    \begin{columns}
        \begin{column}{0.5\textwidth}
            \textbf{Lo que NO haremos:}
            \begin{itemize}
                \item Teoría de espacios vectoriales abstractos.
                \item Inversión de matrices $5\times 5$ a mano.
            \end{itemize}
        \end{column}
        \begin{column}{0.5\textwidth}
            \textbf{Lo que SÍ haremos:}
            \begin{itemize}
                \item \textbf{Mecánica matricial operativa}.
                \item Estructurar sistemas canónicos $\mathbf{A}\mathbf{x} = \mathbf{b}$.
            \end{itemize}
        \end{column}
    \end{columns}
    
    \vspace{0.5cm}
    Del cálculo pesado se encargarán \textbf{Python} y \textbf{LTSpice}. Su trabajo como ingenieros es \textbf{modelar la física correctamente} extrayendo la topología del circuito.
\end{frame}

% ---------------------------------------------------
% Ejercicio Diagnóstico (Dinámica Activa)
% ---------------------------------------------------
\begin{frame}{Desafío Diagnóstico: El Sensor de Dos Nodos (15 min)}
    \begin{columns}
        \begin{column}{0.4\textwidth}
            \begin{center}
            \begin{circuitikz}[scale=0.8, transform shape]
                \draw 
                (0,0) to[V, l=$12\text{V}$] (0,3)
                to[R, l=$R_1$, a=$100\Omega$] (3,3)
                to[short, -*] (3,3) node[above]{$V_1$}
                to[R, l=$R_2$, a=$220\Omega$] (3,0)
                to[short, -*] (3,0)
                (3,3) to[short] (6,3)
                to[R, l=$R_3$, a=$330\Omega$] (6,0)
                (0,0) -- (6,0)
                (3,0) node[ground]{};
            \end{circuitikz}
            \end{center}
        \end{column}
        \begin{column}{0.6\textwidth}
            \textbf{Consigna (En parejas):}
            \begin{enumerate}
                \item Escriban la ecuación de suma de corrientes (LCK) en el nodo $V_1$ ($\sum I = 0$).
                \item Reordenen la ecuación dejando $V_1$ a la izquierda y los términos constantes a la derecha.
                \item Expresen el resultado como: \\ 
                $V_1 \cdot (\text{Conductancia Total}) = I_{\text{fuente}}$
            \end{enumerate}
        \end{column}
    \end{columns}
\end{frame}

% ---------------------------------------------------
% Solución del Ejercicio
% ---------------------------------------------------
\begin{frame}{Construcción Intuitiva de Matrices}
    \textbf{Paso 1:} $\frac{V_1 - 12}{100} + \frac{V_1}{220} + \frac{V_1}{330} = 0$
    
    \vspace{0.3cm}
    \textbf{Paso 2:} $V_1 \left( \frac{1}{100} + \frac{1}{220} + \frac{1}{330} \right) = \frac{12}{100}$
    
    \vspace{0.3cm}
    \textbf{Conclusión:}
    \begin{equation*}
        \mathbf{G}\mathbf{V} = \mathbf{I}
    \end{equation*}
    La matriz de conductancia ($\mathbf{G}$) se forma \textbf{sumando los inversos} de las resistencias conectadas al nodo. ¡Desmitifiquemos las matrices desde el primer día!
\end{frame}

% ---------------------------------------------------
% Tareas y Compromisos
% ---------------------------------------------------
\begin{frame}{Compromisos para la Próxima Sesión}
    \begin{block}{1. Módulo Cero Asincrónico}
        Completar en su tiempo de estudio independiente el \textbf{Jupyter Notebook 1 y 2} (Nivelación en Python y Sistemas Lineales) disponibles en el repositorio.
    \end{block}
    
    \begin{block}{2. Infraestructura Tecnológica}
        Instalar en sus equipos personales:
        \begin{itemize}
            \item \textbf{Python 3.12} (Distribución Anaconda o VS Code).
            \item \textbf{LTSpice XVII}.
        \end{itemize}
    \end{block}
    
    \begin{block}{3. Lectura Previa Fundamental}
        \textbf{C.K. Alexander \& M.N.O. Sadiku.} \textit{Fundamentos de Circuitos Eléctricos}. Leer detenidamente los Capítulos 1 y 2.
    \end{block}
\end{frame}

\end{document}